% ============================================================
% Capítulo: Implementación (Bloque B.9)
% Todos los listados de código de este capítulo son código REAL del
% repositorio (backend-springboot/src/main/java, db/procs/), citados con
% su ruta exacta -- ninguno fue inventado ni parafraseado.
% ============================================================
\chapter{Implementación}
\label{cap:implementacion}

Este capítulo documenta decisiones críticas de implementación del
backend, cada una con al menos un listado de código real extraído
directamente del repositorio.

\section{Procedimiento almacenado y su invocación desde Spring Data JPA}
\label{sec:sp-invocacion}

La estrategia híbrida de acceso a datos (fundamentada en la
\autoref{sec:mt-acceso-datos}; ADR-006, \autoref{sec:adrs})
exige que las operaciones con validación cruzada multi-tabla se
implementen como procedimiento o función SQL. El siguiente listado es
\texttt{db/procs/sp\_crear\_prestamo.sql} completo: registra un préstamo
de forma atómica, validando que el usuario no esté bloqueado por multas
y que el libro tenga stock disponible, usando \texttt{SQLSTATE}
personalizados (\texttt{LB404}/\texttt{LB422}) para que el backend
traduzca el error sin parsear texto (ver Listado \ref{lst:sp-crear-prestamo}).

\begin{lstlisting}[language=SQL, caption={\texttt{db/procs/sp\_crear\_prestamo.sql} (completo)}, label={lst:sp-crear-prestamo}]
CREATE OR REPLACE FUNCTION sp_crear_prestamo(
    p_usuario_id BIGINT,
    p_libro_id BIGINT,
    p_bibliotecario_id BIGINT,
    p_dias_prestamo INTEGER
)
RETURNS BIGINT
LANGUAGE plpgsql
AS $$
DECLARE
    v_estado_usuario_id     INTEGER;
    v_estado_bloqueado_id   INTEGER;
    v_stock_disponible      SMALLINT;
    v_estado_activo_prestamo_id INTEGER;
    v_prestamo_id           BIGINT;
BEGIN
    SELECT id INTO v_estado_bloqueado_id
      FROM estados_usuario
     WHERE nombre = 'BLOQUEADO_POR_MULTA';

    SELECT estado_id INTO v_estado_usuario_id
      FROM usuarios
     WHERE id = p_usuario_id
     FOR UPDATE;

    IF NOT FOUND THEN
        RAISE EXCEPTION 'El usuario % no existe', p_usuario_id
            USING ERRCODE = 'LB404';
    END IF;

    IF v_estado_usuario_id = v_estado_bloqueado_id THEN
        RAISE EXCEPTION 'El usuario % esta bloqueado por multas pendientes ...',
            p_usuario_id USING ERRCODE = 'LB422';
    END IF;

    SELECT stock_disponible INTO v_stock_disponible
      FROM libros WHERE id = p_libro_id FOR UPDATE;

    IF NOT FOUND THEN
        RAISE EXCEPTION 'El libro % no existe', p_libro_id
            USING ERRCODE = 'LB404';
    END IF;

    IF v_stock_disponible <= 0 THEN
        RAISE EXCEPTION 'El libro % no tiene stock disponible', p_libro_id
            USING ERRCODE = 'LB422';
    END IF;

    UPDATE libros SET stock_disponible = stock_disponible - 1
     WHERE id = p_libro_id;

    INSERT INTO prestamos (usuario_id, libro_id, bibliotecario_id,
        fecha_prestamo, fecha_devolucion_estimada, estado_prestamo_id)
    VALUES (p_usuario_id, p_libro_id, p_bibliotecario_id, NOW(),
        NOW() + make_interval(days => p_dias_prestamo),
        v_estado_activo_prestamo_id)
    RETURNING id INTO v_prestamo_id;

    RETURN v_prestamo_id;
END;
$$;
\end{lstlisting}

\textbf{Nota de honestidad sobre el mecanismo de invocación.} La guía de
esta entrega espera que el backend conecte los procedimientos con
\texttt{@Procedure} o \texttt{@NamedStoredProcedureQuery}. Ese es el
contrato vigente para las rutinas principales con efectos secundarios:
\texttt{sp\_crear\_prestamo} y
\texttt{sp\_expirar\_reservaciones\_vencidas} se exponen con
\texttt{@Procedure}; \texttt{sp\_registrar\_devolucion},
\texttt{sp\_pagar\_multa} y \texttt{sp\_anular\_multa} se declaran como
\texttt{@NamedStoredProcedureQuery} en las entidades \texttt{Loan} y
\texttt{Fine}, y sus repositorios los referencian con
\texttt{@Procedure}. El siguiente listado muestra el método vigente hoy en
\texttt{LoanProcedureRepository.java}:

\begin{lstlisting}[language=Java, caption={\texttt{LoanProcedureRepository.spCreateLoanProcedure} (mecanismo vigente)}, label={lst:sp-crear-prestamo-java}]
@org.springframework.stereotype.Repository
public interface LoanProcedureRepository
        extends Repository<Loan, Long>, LoanProcedureRepositoryCustom {

    @Procedure(procedureName = "sp_crear_prestamo")
    Long spCreateLoanProcedure(Long userId, Long bookId,
                               Long librarianId, Integer daysLoan);

    @Procedure(name = "Prestamo.registrarDevolucion")
    Map<String, Object> spRegisterLoanReturn(Long loanId);
}
\end{lstlisting}

Las funciones de reporte/listado que retornan \texttt{TABLE} (varias
filas) permanecen con \texttt{@Query(nativeQuery = true)} por una
limitación distinta: JPA~2.1 no expone de forma estándar un
\texttt{RETURNS TABLE} de PostgreSQL mediante \texttt{@Procedure} salvo
reescribiendo esas funciones a \texttt{REF\_CURSOR}. Esa excepción queda
acotada a rutinas tabulares de lectura, no a las operaciones principales
con efectos secundarios.

\subsection*{¿Esto viola la prohibición de SQL dinámico de la guía (A.2.1)?}

No. La guía (Bloque~A.2.1) establece textualmente que ``queda prohibido
invocarlos mediante concatenación de cadenas en
\texttt{createNativeQuery(...)}''. Esa prohibición apunta a un patrón
concreto y distinto del usado en este proyecto: construir la sentencia
SQL en tiempo de ejecución \textbf{concatenando texto de entrada del
usuario} dentro del propio literal de la consulta -- p.~ej.
\texttt{createNativeQuery("SELECT sp\_crear\_prestamo(" + idUsuario +
")")}, donde \texttt{idUsuario} pasa a formar parte literal de la
cadena SQL antes de ejecutarse, el vector clásico de inyección SQL.

El código real de este proyecto \textbf{no hace eso en ningún punto}. En
las rutinas principales se usa \texttt{@Procedure}/
\texttt{@NamedStoredProcedureQuery}; en las funciones tabulares que usan
\texttt{@Query(nativeQuery = true)}, el literal SQL es fijo y los
parámetros se bindean por separado. El valor de cada parámetro nunca se
concatena, formatea ni interpola dentro de la cadena SQL.

En síntesis: \textbf{la regla A.2.1 prohíbe construir SQL concatenando
entrada de usuario dentro del texto de la consulta}. Este proyecto
mantiene el mecanismo exigido para las rutinas principales y conserva
\texttt{@Query} solo donde la forma tabular de PostgreSQL lo justifica.
El CI ejecuta \texttt{scripts/audit-sql-dynamic.sh} para rechazar SQL
dinámico en \texttt{db/procs/}.

\section{Manejo uniforme de errores (RFC 7807)}

Todas las respuestas de error del backend se devuelven como
\texttt{ProblemDetail} (RFC~7807), vía un único
\texttt{@RestControllerAdvice} (\texttt{GlobalExceptionHandler}). Un
caso particularmente representativo es la traducción de los
\texttt{SQLSTATE} personalizados de los procedimientos almacenados
(\autoref{sec:sp-invocacion}) a códigos HTTP, sin que el backend tenga
que parsear el texto del mensaje de error (ver Listado \ref{lst:global-exception-handler}):

\begin{lstlisting}[language=Java, caption={\texttt{GlobalExceptionHandler.handleStoredProcedureError} (extracto real)}, label={lst:global-exception-handler}]
@ExceptionHandler({
        InvalidDataAccessResourceUsageException.class,
        UncategorizedSQLException.class,
        DataAccessException.class
})
public ProblemDetail handleStoredProcedureError(Exception ex) {
    Throwable causa = ex;
    while (causa != null && !(causa instanceof SQLException)) {
        causa = causa.getCause();
    }

    if (causa instanceof SQLException sqlEx) {
        String sqlState = sqlEx.getSQLState();
        HttpStatus status = switch (sqlState) {
            case "LB404" -> HttpStatus.NOT_FOUND;
            case "LB409" -> HttpStatus.CONFLICT;
            case "LB422" -> HttpStatus.UNPROCESSABLE_ENTITY;
            default -> null;
        };
        if (status != null) {
            return ProblemDetail.forStatusAndDetail(status, sqlEx.getMessage());
        }
    }

    log.error("Error no controlado en procedimiento almacenado", ex);
    return ProblemDetail.forStatusAndDetail(
            HttpStatus.INTERNAL_SERVER_ERROR, "Error interno del servidor");
}
\end{lstlisting}

Sin este manejador dedicado, cualquier \texttt{RAISE EXCEPTION ... USING
ERRCODE = 'LBxxx'} de un procedimiento llegaba envuelto en una
\texttt{DataAccessException} genérica de Spring, y caía en el
\emph{catch-all} de último recurso -- un error de negocio real (p.~ej.
``préstamo ya devuelto'') se veía como un \texttt{500} genérico en vez
de un \texttt{409} con el mensaje real. El mismo patrón de un
\texttt{@ExceptionHandler} dedicado por tipo de excepción se repite para
excepciones de dominio Java (p.~ej.
\texttt{CorreoYaRegistradoException} $\to$ \texttt{409},
\texttt{LoginRateLimitExcedidoException} $\to$ \texttt{429}) y para
excepciones del propio framework que antes caían también en el
\emph{catch-all}: \texttt{AuthorizationDeniedException} (autorización de
método vía \texttt{@PreAuthorize}, antes producía \texttt{500} en vez de
\texttt{403}) y \texttt{NoResourceFoundException} (recurso estático
inexistente bajo el perfil \texttt{prod}, corregido en esta misma
entrega -- ver \texttt{docs/mediciones/sec/2026-08-11-owasp-a05-verificacion-real.md}).

\section{Configuración paramétrica del sistema}

\texttt{ConfiguracionSistemaService} expone parámetros de la tabla
\texttt{configuracion\_sistema} (p.~ej. el máximo de renovaciones de un
préstamo, ver \texttt{REQ-F-018}) para leerse en cada operación de
negocio sin ir a la base de datos cada vez, con un caché en memoria
simple (\texttt{ConcurrentHashMap} por instancia, no Redis) invalidado
clave por clave al escribir (ver Listado \ref{lst:config-service}):

\begin{lstlisting}[language=Java, caption={\texttt{ConfiguracionSistemaService} (extracto real: lectura cacheada y escritura)}, label={lst:config-service}]
private final ConcurrentHashMap<String, String> cache = new ConcurrentHashMap<>();

@Transactional(readOnly = true)
public String obtenerValor(String clave) {
    String cacheado = cache.get(clave);
    if (cacheado != null) {
        return cacheado;
    }
    String valor = repo.findById(clave)
            .orElseThrow(() -> new EntityNotFoundException(CLAVE_NO_ENCONTRADA + clave))
            .getValor();
    cache.put(clave, valor);
    return valor;
}

@Transactional
public ConfiguracionSistemaResponseDTO actualizar(String clave, String nuevoValor) {
    ConfiguracionSistema config = repo.findById(clave)
            .orElseThrow(() -> new EntityNotFoundException(CLAVE_NO_ENCONTRADA + clave));
    config.setValor(nuevoValor);
    repo.save(config);
    cache.remove(clave);
    return new ConfiguracionSistemaResponseDTO(config.getClave(), config.getValor());
}
\end{lstlisting}

La elección deliberada de un \texttt{ConcurrentHashMap} local en vez de
Redis (que sí se usa para el caché del catálogo, ver
\autoref{sec:cache-redis}) está documentada en el propio Javadoc de la
clase: este valor cambia con muy poca frecuencia (un \texttt{ADMIN}
ajustando un parámetro puntual), a diferencia del caché ``libros'', que
sí necesita compartirse entre instancias e invalidarse por TTL. Además,
\texttt{actualizar} solo permite modificar claves \textbf{ya
existentes} -- el catálogo de parámetros válidos lo define el dominio
(migraciones/\emph{seed}), no quien llama al endpoint, evitando que
quede una clave suelta sin ningún consumidor real.

\section{Estrategia de caché Redis del catálogo}
\label{sec:cache-redis}

El listado de libros usa caché declarativo de Spring
(\texttt{@Cacheable}/\texttt{@CacheEvict}) sobre Redis, con TTL
configurable externamente (ADR-008, ver Listado \ref{lst:libro-service-cache}):

\begin{lstlisting}[language=Java, caption={\texttt{LibroService} (extracto real: lectura cacheada y evicción en escritura)}, label={lst:libro-service-cache}]
@Cacheable("libros")
@Transactional(readOnly = true)
public Page<LibroResponseDTO> listar(Pageable pageable) {
    return libroRepo.findByEstado_Nombre(ESTADO_ACTIVO, pageable)
            .map(this::toDTO);
}

@CacheEvict(value = "libros", allEntries = true)
@Transactional
public LibroResponseDTO crear(LibroRequestDTO dto) {
    if (libroRepo.existsByIsbn(dto.isbn())) {
        throw new IllegalArgumentException("ISBN ya registrado: " + dto.isbn());
    }
    validarStock(dto.stockTotal(), dto.stockDisponible());
    return toDTO(libroRepo.save(fromDTO(dto)));
}
\end{lstlisting}

\texttt{crear}, \texttt{actualizar} y \texttt{eliminar} invalidan el
caché completo del listado (\texttt{allEntries = true}) para no dejar
una respuesta paginada obsoleta tras cualquier escritura. El método
\texttt{listarPorCategoria} (filtro por categoría/autor) deliberadamente
\textbf{no} lleva \texttt{@Cacheable}: combinar el filtro con el caché
exigiría una clave compuesta fuera del alcance de la rama que lo
introdujo, y no es el camino más transitado del catálogo -- decisión
documentada como comentario en el propio código, no un descuido. El
método \texttt{sugerir} (autocompletado) usa un caché \emph{distinto}
(\texttt{"sugerencias-libros"}) con TTL propio, mucho más corto (5--10\,s),
porque su patrón de invalidación (una entrada por texto de búsqueda) es
incompatible con el caché de listado general.

\section{Transporte del token de sesión en cookie \texttt{HttpOnly}}

El \texttt{refreshToken} viaja exclusivamente en una cookie
\texttt{HttpOnly}, \texttt{Secure}, \texttt{SameSite=None}, con
\texttt{path=/api/auth} (ADR-012, ver Listado \ref{lst:jwt-cookie}) -- nunca en el cuerpo JSON, para que
JavaScript del frontend no pueda leerlo ni reenviarlo manualmente:

\begin{lstlisting}[language=Java, caption={\texttt{AuthController.buildRefreshCookie} (real, completo)}, label={lst:jwt-cookie}]
private ResponseCookie buildRefreshCookie(String valor, long maxAgeMs) {
    return ResponseCookie.from(REFRESH_COOKIE_NAME, valor)
            .httpOnly(true)
            .secure(true)
            .sameSite("None")
            .path("/api/auth")
            .maxAge(Duration.ofMillis(maxAgeMs))
            .build();
}
\end{lstlisting}

El mismo método se reutiliza para \emph{limpiar} la cookie en
\texttt{logout} (\texttt{buildRefreshCookie("", 0)}, valor vacío y
\texttt{maxAge=0}). El \texttt{accessToken}, de vida corta (1\,h), sigue
viajando en el cuerpo JSON de la respuesta -- migrarlo también a cookie
está deliberadamente diferido (ver la nota de honestidad de
\texttt{ADR-012} y \texttt{REQ-NF-002}, \autoref{cap:ingenieria-requisitos})
por el impacto que tendría en \texttt{jwt.interceptor.ts}/
\texttt{auth.service.ts} del frontend, no por un olvido.
